\section{Preliminaries}
\label{sec:prelim}

\subsection{Notation}

\begin{table}[h]
\centering
\begin{tabular}{ll}
\toprule
Symbol & Meaning \\
\midrule
$\Omega \subset \R^d$ & spatial domain (bounded, periodic, or torus) \\
$T > 0$               & finite time horizon \\
$\rho(t,x)$           & population density at $(t,x)$, with $\int_\Omega \rho(t,\cdot) = 1$ \\
$\varphi(t,x)$        & value function (HJB unknown) \\
$\nu \geq 0$          & diffusion coefficient \\
$H(x,\rho,p)$         & Hamiltonian; argument $p$ is the momentum $\nabla\varphi$ \\
$H_0(x,p)$            & separable base Hamiltonian \\
$f_0(x,\rho)$         & monotone congestion interaction \\
$R(x,\rho,p)$         & coupling residual \\
$\eps \geq 0$         & coupling strength \\
$c_0 > 0$             & strong-convexity constant of $H_0$ in $p$ \\
$\lambda > 0$         & Lasry-Lions monotonicity constant of $f_0$ \\
$L_R, L_{R,p}, L_{p\rho} \geq 0$ & constants controlling $R$, its $p$-derivatives, and $\nabla_p R$ \\
$\rho_{\min}, \rho_{\max}$ & lower/upper density bounds (non-vacuum condition) \\
$\kappa_\eps$ & contraction factor of the exact outer iteration \\
$\mathcal{T}$         & outer fixed-point map  $\bar\rho \mapsto \rho$ \\
$\mathcal{W}_2$       & Wasserstein-2 distance \\
\bottomrule
\end{tabular}
\caption{Notation used throughout the paper.}
\label{tab:notation}
\end{table}

\subsection{Mean Field Games}

Let $\Omega \subset \R^d$ be a bounded domain and $T > 0$ a finite horizon.
A Mean Field Game is described by the coupled PDE system
\eqref{eq:HJB}--\eqref{eq:FP} with initial condition
$\rho(0,\cdot) = \rho_0$ and terminal condition
$\varphi(T,\cdot) = g(\cdot)$.
The solution $(\varphi, \rho)$ represents a Nash equilibrium: no agent
can improve their cost by unilaterally changing their control.
Existence and uniqueness under monotonicity conditions on the coupling
are established in \citet{lasry2007mean}.  Density-dependent terminal costs
can be included under an additional terminal monotonicity condition; we keep
$g$ fixed to match the stability argument used below.

\subsection{Separable Hamiltonians and GAN-Based Methods}

When the Hamiltonian has separable structure
$H(x,\rho,p) = H_0(x,p) - f_0(x,\rho)$,
the MFG system admits a primal-dual variational formulation.
Specifically, \citet{cao2021connecting} show that the MFG can be written as
\begin{equation}\label{eq:saddle}
  \inf_{\rho}\, \sup_{\varphi}\;
  \Bigl\{
    \E_{z\sim P(z),\, t\sim\mathrm{Unif}[0,T]}
    \bigl[\partial_t\varphi + \nu\Delta\varphi - H_0(\cdot,\nabla\varphi)\bigr]
    + \E_z[\varphi(0,\cdot)] - \E_{x\sim\rho_T}[\varphi(T,x)]
    + \E[f_0(\cdot,\rho)]
  \Bigr\}.
\end{equation}
The min-max structure of \eqref{eq:saddle} matches the
Kantorovich-Rubinstein dual formulation of Wasserstein GANs,
enabling two families of GAN-based solvers.

\paragraph{APAC-Net \citep{lin2021apac}.}
Uses a generator network to parameterize $\rho$ and a critic to parameterize
$\varphi$, trained adversarially to solve \eqref{eq:saddle}.
Convergence is established in the Wasserstein-1 metric.

\paragraph{MFGAN \citep{ruthotto2020machine}.}
A related adversarial framework that exploits the connection to
Wasserstein-2 optimal transport for mean field control problems.

\paragraph{Hard limitation.}
Both methods require separability of $H$.
For non-separable $H$, the coupling between $\rho$ and $p$ prevents
the saddle-point formulation \eqref{eq:saddle} from being derived,
since the Hopf-type formula used in \citet{cao2021connecting}
is only valid when $H$ splits as $H_0(x,p) - f_0(x,\rho)$.

\subsection{Non-Separable Hamiltonians: MFDGM}

\citet{assouline2023deep} propose MFDGM, which uses two neural networks
$N_\omega$ and $N_\theta$ to approximate $(\varphi, \rho)$ by minimizing
the PDE residuals directly:
\begin{align*}
  \mathcal{L}_{\mathrm{HJB}} &= \frac{1}{B}\sum_{b=1}^B
    \bigl|\partial_t\varphi_\omega + \nu\Delta\varphi_\omega
    - H(x_b,\rho_\theta,\nabla\varphi_\omega)\bigr|^2
    + \frac{1}{S}\sum_{s=1}^S
    \bigl|\varphi_\omega(T,x_s) - g(x_s)\bigr|^2,\\
  \mathcal{L}_{\mathrm{FP}} &= \frac{1}{B}\sum_{b=1}^B
    \bigl|\partial_t\rho_\theta - \nu\Delta\rho_\theta
    - \Div(\rho_\theta\,\nabla_p H)\bigr|^2
    + \frac{1}{S}\sum_{s=1}^S
    \bigl|\rho_\theta(0,x_s) - \rho_0(x_s)\bigr|^2.
\end{align*}
MFDGM handles any $H$, separable or not, but at a cost:
it approximates $\rho$ pointwise rather than distributionally,
and its convergence (Theorem~3.2 of \citealt{assouline2023deep})
holds only in $L^p$ for $p < 2$, with no explicit rate.

\subsection{Summary: The Gap Between the Two Families}

\begin{table}[h]
\centering
\begin{tabular}{lcccc}
\toprule
Method & Non-sep.\ $H$ & $\rho$ representation & Convergence metric & Rate \\
\midrule
APAC-Net / MFGAN & No  & Distributional & Wasserstein-1 & Implicit \\
MFDGM            & Yes & Pointwise NN   & $L^p$, $p<2$  & None \\
Ours & Weak & Compatible with distributional inner solvers & Wasserstein-2 outer map & Explicit \\
\bottomrule
\end{tabular}
\caption{Comparison of deep learning solvers for MFGs.
``Weak'' refers to the GAN-tractable regime defined in \Cref{sec:weak}.}
\label{tab:comparison}
\end{table}

\noindent The core gap is that neither existing family simultaneously handles
non-separable Hamiltonians and retains the distributional structure of
adversarial MFG solvers. Our contribution addresses the weakly
non-separable regime by proving an explicit contraction rate for the
outer reduction; the distributional representation is inherited whenever
the frozen separable subproblems are solved by a GAN-based inner method.
